\chapter{“二十秒”}
\label{sec:1}
一名情报军官在讲述自己在一条以机器生成的猎杀名单为核心搭建起来的目标选定流水线中的角色时，对记者说，他在每个名字上花二十秒，刚好够他确认目标是个男人：“要分辨一个人是男声还是女声，用不了多长时间。”他还说，作为一个人，他的附加价值是零，除了当一枚批准图章之外。“这节省了大量时间”\autocite{abraham2024lavender}。按他的说法，在作战节奏之下，“人在回路中”实际上就是这么回事。\footnote{关于自动化偏差的研究预测了这种委托：当其他任务在争夺人的注意力时，拿到机器建议的人会更少去核对它，也更听从它\autocite{parasuraman2010complacency}。}

一种解读是，这名军官是一个被给了太少时间的主体，解决办法是把他管理得更好——更多时间、第二名审查员、在回路中加装更多良知。通过这种\emph{社会屈从}的视角，人是一个站在机器之外、把机器当作工具来操作的主体：司机与汽车，工人与车床。这个人有角色、有名字、有职责，还有一个可以被唤起的良知。另一种视角是\emph{机器奴役}：人是机器的一个组件，和机器的其他零件组装在一起，成为一个把它们一并调节的系统；机器把人当作自己的器官之一来使用，而人所贡献的，是一项按系统设定的速率执行的功能。\footnote{多克托罗给这名军官的位置起的名字是“反向半人马”：不是由机器辅助的人，而是由人做成的机器组件\autocite{doctorow2026reversecentaur}。}两者并不互相排斥；一个人可以在一个层面上担任一个角色，同时在另一个层面上是一个零件\autocite[高原13，“The Apparatus of Capture”：当人类本身就是一台机器的组成部件，这台机器由他们彼此之间以及与其他事物共同构成，并受一个更高的统一体控制和指挥时，就有了奴役；当这个更高的统一体把人构建为一个主体，与一个新的外部客体相连——这个客体可以是一只动物、一件工具，甚至是一台机器——时，就有了屈从]{deleuze1987plateaus}，而那名二十秒军官两者兼而有之。作为零件，他就是流水线盖下批准图章的那个位置，而二十秒就是这个零件的周期时间。作为主体，他自始至终都被唤起：他向一名记者讲述了平民死亡的消息与工作中要求更多目标的压力之间的落差。而即便这种落差触及了他的良知，他还是照做了。它被感受到了，却什么也没有决定；\footnote{或者至少，没有决定任何能阻止以色列国防军在2023年10月至2025年10月间在加沙杀死10万人的事\autocite{gomezugarte2025uncertainty,acosta2026gaza,economist2025gaza}。而那名军官的讲述，是+972 Magazine在2024年4月发表的。他决定向记者开口，到底让死亡人数少了一个、少了5万个，还是一个也没少，谁也说不准。而如果那名军官在任职两周后就拒绝——比如说，假设是在2023年11月——又会怎样？我们说的是一名用Lavender来汇编猎杀名单的情报军官。他日常的工作环境，大概更像一份坐办公室的文职工作，而不是动作片。他是否应该沿着指挥链一路往上，一再要求与上级面谈，同时拒绝批准那些可以合理预见会造成平民死亡的空袭目标？可以说，那才是合乎伦理的做法。在那个反事实情境里，死的人会更少吗？}他说，认定那项政策不成比例的判断，是事后回过头来才有的。

所以缺的并不是良知；他有良知。良知起作用的方式，是让一个决定付出代价，但在这里，代价的净值指向了另一边：“我们一直在被施压：‘给我们更多目标。’他们真的冲我们大吼大叫”\autocite{abraham2024lavender}。他是一个有顾虑的组件，但节奏保住了：那些顾虑丝毫没有改变他所贡献的东西，于是他成了一枚橡皮图章。一个有顾虑、却跟不上节奏的组件，就会变成以色列信号情报部队一名前指挥官写了一本书、主张要去除的那个“人类瓶颈”\autocite{ys2021humanmachine}。让这名军官成为组件的，是那套安排，而不是他的态度。组件可以不情愿地照做，也可以热切地照做：阿伦特在艾希曼审判中看到的那个不思考的公务员\autocite{arendt1963eichmann}，从他四年前在阿根廷、在毫无利害关系的情况下所说的话来看，是他为耶路撒冷演的一出戏\autocite{stangneth2014eichmann}。流水线让这名军官能贡献的，是一个按固定速率执行的固定操作：他看到的每一个案件，都是流水线所描述的样子，一次一件，在流水线给他的时间之内。是否有任何东西能站在那个位置上而不成为组件，是第四部要讨论的问题。

调查报道描述了在加沙协同运作的三套系统：一套生成建筑物和基础设施目标，据消息来源称，其中大多数建议是私人住宅；第二套按估算的武装组织关联程度给个人打分，据说它生成的名单标记了多达3.7万人，内部的认识是大约每十次错一次\footnote{以色列军方对这一说法提出异议。下文关于米纳卜和Maven的叙述，所依据的调查报道同样遭到当事各方的质疑。我按报道所述采用。}；第三套则追踪被标记的人什么时候回家，好让空袭能选在他们和家人在一起的时候进行。

\runin{Maven：美国版，现在加上了生成式AI加速}

米纳卜那所学校遭到的轰炸造成156人死亡，其中120人是学童\autocite[第92段]{UNFFMIran2026}。\footnote{数字因来源和时间而异，从156人（米纳卜检察官给出的最终死亡人数，国际特赦组织也采用这一数字）到168人（该镇镇长在最初几天给出的数字）。联合国调查团采用检察官的数字，同时也报告了独立信息显示有157人死亡，其中123人是13岁及以下的儿童\autocite[第92–93段]{UNFFMIran2026}。}这所学校的民用功能在开源卫星图像中清晰可见，而目标编码仍带着一个从相邻的革命卫队设施沿袭下来的标识。联合国一个实况调查团报告称，有合理理由相信这次轰炸构成以民用物体为攻击对象的战争罪\autocite{ohchr2026minabpress}，并指出这所学校清晰可辨，没有任何信息表明它被用于军事用途\autocite[第81–86、124–126段]{UNFFMIran2026}。调查团认定，美国未能履行尽一切可行措施核实目标的义务，并把更新目标数据列为可行的预防措施之一，还认为这一失职已超出疏忽的程度\autocite[第124–125段]{UNFFMIran2026}。国际特赦组织得出了类似的结论，其措辞与《第一附加议定书》第57条第2款(a)项(i)目规定的核实义务相呼应：部队可能依赖了过时的情报，且未能尽一切可行措施核实目标\autocite{Amnesty2026Minab}。

事情发生在2月28日上午，伊朗战役的第一天\autocite{Amnesty2026Minab,crs2026iranwar}。前一天，国防部长宣布，那家供应商（Anthropic）将被认定为供应链风险，理由是它拒绝许可Claude用于一切合法用途；这家供应商的模型（Claude）就在那个目标选定平台（Maven Smart System）里面\autocite{npr2026openai}。Maven Smart System曾把这所学校列为建议的首日目标\autocite{livingstone2026maven,gizmodo2026minab}。没有任何已发表的报道说明Claude在这项建议中起了什么作用，而一位熟悉该供应商工作的人士告诉NBC News，该模型并没有直接提供目标选定建议\autocite{nbc2026targeting}。

一项五角大楼内部审查——其完整结论至今仍被严格保密——据报认定：在以往的冲突中需要几个小时的目标准备工作，被压缩到了几分钟；平民伤害缓解人员编制减少了大约九成，降到不足20人；中央司令部的团队从10人减到1人；而且团队中没有任何人审查过米纳卜这个地点——官员把这一结果既归咎于人员削减，也归咎于一项不让该团队参与计划和作战的决定\autocite{Bloomberg2026KillChain}。正式调查已于3月启动，此前一项初步评估认定美国很可能负有责任\autocite{NPR2026MinabProbe}。

在2026年的伊朗战役中，一个由承包商搭建、内含一家商业供应商的通用模型的目标选定平台，被用来按战略重要性给目标排序，并为每一次打击起草法律依据\autocite{livingstone2026maven}；它帮助在最初24小时内打击了一千多个目标，大约是它出现之前所能做到的十倍\autocite{csis2026maven}，到第38天时已达1.3万个\autocite{breakingdefense2026maven}。五角大楼没有公布它对每个目标的审查时间。\footnote{Maven是一个承载着不止一套系统的项目名称：一个是至今仍叫Maven的五角大楼项目，另一个是从中发展出来的软件平台Maven Smart System\autocite{csis2026maven}。这个项目始于2017年，最初是一个在无人机于打击伊斯兰国的行动中所采集的视频里标注物体的项目\autocite{weisgerber2017maven,fang2019maven}。2026年的系统把各种军事系统和情报来源融合进同一个界面，用来寻找并打击目标以及指挥作战\autocite{defensescoop2026transition}。这个平台是Palantir的：五角大楼的首席数字与人工智能官称它为“Palantir的软件即服务产品”\autocite{defensescoop2026transition}。Palantir的公开立场（经转述）是，决定生死不是它的职责，责任始终在军事机构身上\autocite{lee2026palantir}；它还坚称底层情报是政府的责任，并且没有证据表明Maven出了故障或导致了米纳卜的轰炸\autocite{Bloomberg2026KillChain}。这个项目如何从一个办公室和承包商转到下一个，见附录~\ref{sec:C.1}。}

以色列和美国空袭名单上的人，不可能对自己上了名单提出异议——那些流水线本来就不是为了获取这种反馈而建的。安德里·何塞·埃尔南德斯·罗梅罗（Andry José Hernández Romero）被卷进的是另一条流水线，而他本来是可以提出异议的。他是来自委内瑞拉的化妆师，2024年8月前往圣伊西德罗入境口岸申请庇护，通过了可信恐惧面谈，有律师，有待审的申请，也有听证日期\autocite{LULAC2025HernandezRomero,Vega2025HernandezRomero}。一名负责对他进行筛查的承包商雇员，把他的两处文身——“妈妈”和“爸爸”字样上方的王冠——解读为帮派标记\autocite{Palomera2025Blade}；2025年3月，他在听证之前就被送上飞机，飞往萨尔瓦多的一座监狱，就在一名联邦法官下令让飞机掉头的那个晚上\autocite{jgg2025contempt}。一个人要提出异议所需要的一切，他都有，除了时间。等我们讨论美国的绑架与强制移送流水线时，会更完整地回到他的案子上来。

\textbf{谁拒绝了。}二十秒的记录还有另一半，也就是谁顶住了那种压力的记录，而它告诉我们，房间的构成非常重要。在米尔格拉姆的基线条件下，大约三分之一的被试在电击达到最大值之前就中止了。当有两名同伴抗命时，只有十分之一的人坚持到最后；而当两名实验者意见不一时，继续下去的人远比基线条件下少\autocite{milgram1974obedience}。阿希的被试中，大约四分之一从未向错误的多数屈服\autocite{asch1956studies}。布朗宁估计，第101预备警察营的成员中，有十分之一到五分之一的人拒绝或逃避了在约瑟夫乌（Józefów）的枪杀，那是该营的第一场屠杀\autocite{browning1992ordinary}。这个比例随情境而变，从几乎没有到几乎全部。那些顶住了的人，手里有来自渠道之外的证据：一个抗命的同伴，一个意见不同的第二权威，以及——在奥利纳夫妇研究的救助者当中——跨越了迫害所划界线的人际关系\autocite{oliner1988altruistic}。或者，他们以压力无法压倒的精确度坚守承诺，就像人们坚守道德使命和神圣价值那样\autocite{skitka2009limits,ginges2007sacred}。或者，他们把信息来源看作有利害关系的一方。而群体在孤立的个人很少能抵抗的地方抵抗了\autocite{gamson1982encounters}。

房间同样可以为另一方布置。魏玛德国有着稠密的结社生活，而这或许助推了共和国的垮台：当政党和政府回应不了人们的不满，人们就把精力投进按群体界线组织起来的俱乐部和社团，而纳粹则通过这些社团早已建立起来的网络招募成员\autocite{berman1997weimar}。合唱团、体育俱乐部和体操协会越多的城镇，入党的速度就越快，但只限于政府不稳定的地方；在政府稳定的普鲁士，这种关联很小，也不显著\autocite{satyanath2017bowling}。所以救助者所拥有的，并不只是有更多人做伴。他们的纽带跨越了迫害所划的界线。阿舒托什·瓦尔什尼在印度城市中也发现了同样的情况：在社团把印度教徒和穆斯林两个群体纳入同一些组织的地方，两者之间的骚乱很少发生；在没有这样做的地方，骚乱则很常见\autocite{varshney2002ethnic}。让拒绝成为可能的那些纽带，正是一场法西斯运动必须切断的纽带。

接下来的一切，都是在讲怎样布置房间。\footnote{房间就是场景，而正如我的戏剧老师会说的，布景决定这场戏怎么演。房间里还有谁？离开之后能不能活下去？会有人听见拒绝吗？}“布置”是一个关于制作的词，而汉娜·阿伦特曾告诫，不要把政治当成一件被制作出来的东西。她写道，对希腊人来说，法律字面上就是一道墙，没有它，也许能有一座城镇，却不会有一个城邦\autocite{arendt1958human}：立法者筑起了墙，而墙内发生的事，不是他们能制作的。我所主张的那些基本措施，就是这个意义上的墙。它们让人拒绝之后能活下去，让拒绝被听见；它们不会产生拒绝，而按她的说法，权力只存在于共同行动的人们之间，他们一散，权力就没了——这正是那份谁拒绝了的记录所显示的、关于拒绝的事。第二部和第三部的工具让拒绝起作用。第五部的制度造就拒绝者。一台能够拒绝的机器，需要的是那些身边有同伴的十分之九的人所拥有的东西：身边一个经历不同的同伴，一个可以去的地方，还有一个会听见它的人。

这也正是流水线中的一个拒绝者——包括一台拒绝的机器——对那名军官可能意味着的东西。案件到他手上时，一台拒绝的机器会在屏幕上：那个抗命的同伴——前提是这种效应从房间搬到屏幕上之后依然成立；这一点没有经过检验，而自动化偏差预测它不会成立。一台靠审查员的顺从来左右他们的机器，只会把他们变成方向相反的橡皮图章。那名只有二十秒、身边有一台拒绝的机器的军官，还是报道中的那名军官吗？
